Kita tunjukkan bahwa untuk setiap titik

\mavergleichskette
{\vergleichskette
{ x
}
{ \in }{ U
}
{ = }{ X \setminus Y
}
{    }{
}
{   }{
}
}
{}{}{}
terdapat sebuah lingkungan terbuka

\mavergleichskette
{\vergleichskette
{ x
}
{ \in }{ V
}
{ \subseteq }{ U
}
{    }{
}
{   }{
}
}
{}{}{}.
Gabungan semua lingkungan tersebut kemudian sama dengan $U$. Jadi, misalkan

\mavergleichskette
{\vergleichskette
{ x
}
{ \notin }{ Y
}
{  }{
}
{    }{
}
{   }{
}
}
{}{}{}
diberikan. Untuk setiap titik

\mavergleichskette
{\vergleichskette
{ y
}
{ \in }{ Y
}
{  }{
}
{    }{
}
{   }{
}
}
{}{}{}
terdapat himpunan-himpunan terbuka

\mavergleichskette
{\vergleichskette
{ x
}
{ \in }{ V_y
}
{  }{
}
{    }{
}
{   }{
}
}
{}{}{}
dan

\mavergleichskette
{\vergleichskette
{ y
}
{ \in }{ W_y
}
{  }{
}
{    }{
}
{   }{
}
}
{}{}{,}
yang saling lepas. Himpunan-himpunan terbuka

\mathbed {W_y} {}
 {y \in Y} {}
 {} {} {} {,}
menutupi $Y$. Berdasarkan kekompakan, terdapat sebuah subtutupan hingga, katakanlah

\mavergleichskettedisp
{\vergleichskette
{ Y
}
{ \subseteq} { \bigcup_{i = 1}^n  W_i
}
{ } {
}
{ } {
}
{ } {
}
}
{}{}{.}
Maka

\mavergleichskette
{\vergleichskette
{ V
}
{ = }{ \bigcap_{i = 1}^n  V_i
}
{  }{
}
{    }{
}
{   }{
}
}
{}{}{}
terbuka sebagai irisan berhingga dari himpunan-himpunan terbuka, dan karena itu merupakan sebuah lingkungan terbuka dari $x$ yang saling lepas dengan $Y$.
